% =====================================================================
%  Guide — Environment Setup: Flexus X + SSH Proxy + MaaS Gateway + Huawei Project
%  Compile with XeLaTeX:
%      xelatex setup-guide.tex   (run 2x for the TOC)
% =====================================================================
\documentclass{guide}

% --- Document configuration -----------------------------------------
\setguidetitle{Guide: Environment Setup for Huawei Cloud Development}
\setheadertitle{Huawei Cloud -- Flexus X + SSH + MaaS Gateway Setup}
\setcovertext{Huawei Technologies CO., LTD}
\setdocversion{3.1.1}

\begin{document}

% ---- COVER ---------------------------------------------------------
\makecover

% ---- TABLE OF CONTENTS ---------------------------------------------
\maketoc

% ---- BODY (restarts page numbering at 1) ---------------------------
\startbody

% =====================================================================
%  Chapter 1 — Provision a Flexus X Instance
% =====================================================================
\section{Provision a Flexus X Instance on Huawei Cloud}

\begin{objectives}
  \generalobjective{Provision a Flexus X instance on Huawei Cloud that will
  host a complete coding stack environment: the
  \weblink{https://github.com/wallacelw/oh-my-coding-maas-gateway}{oh-my-coding-maas-gateway}
  (a LiteLLM-based proxy with load balancing, virtual keys, and Grafana
  observability), governance tooling, and multiple coding agents --- with
  OpenCode as the primary one. By the end of this chapter you will have a
  running Flexus X instance with a public IP accessible over SSH.}

  \prerequisites
  \begin{itemize}
    \item Active Huawei Cloud account with IAM credentials.
    \item Permission to create security groups, Flexus X instances, and EIPs.
    \item A key pair created in the Console, or permission to create one.
    \item A web browser (Chrome, Firefox, or Edge).
  \end{itemize}
\end{objectives}

\begin{infobox}
This guide assumes the \inlinecode{sa-brazil-1} region. Adjust the region if your
account is registered elsewhere.
\end{infobox}

\subsection{Create the security group (sg-xgate)}

\objective{Create a security group that allows SSH access on port 4444 from
specific Huawei internal IPs only.}

\begin{infobox}
We use port \inlinecode{4444} instead of the standard SSH port \inlinecode{22} because
the Huawei corporate proxy blocks outbound connections on port 22. Port
\inlinecode{4444} is allowed through the proxy and works reliably for SSH.
\end{infobox}

\stepbystep
\begin{enumerate}
  \item In the Console, go to \menu{VPC, Security Groups}
        and click \textbf{Create Security Group}.
  \item Name it \inlinecode{sg-xgate} and add the following inbound rules:

  \begin{code}
Rule 1:
  Protocol:    TCP
  Port:        4444
  Source:      119.8.89.193/32
  Description: SSH for Huawei internal user 1

Rule 2:
  Protocol:    TCP
  Port:        4444
  Source:      119.8.89.3/32
  Description: SSH for Huawei internal user 2

Rule 3:
  Protocol:    TCP
  Port:        4444
  Source:      119.8.193.3/32
  Description: SSH for Huawei internal user 3
  \end{code}

  \item Click \textbf{OK} to create the security group.
\end{enumerate}

\begin{tip}
Each \inlinecode{/32} rule restricts access to exactly one IP address. Only the
three listed Huawei internal IPs can reach the instance on port 4444 --- no other
inbound traffic is allowed.
\end{tip}

\subsection{Create the Flexus X instance}

\objective{Launch a single Flexus X instance using the Console wizard, using the
security group created above.}

\stepbystep
\begin{enumerate}
  \item Log in to the \weblink{https://console.huaweicloud.com/}{Huawei Cloud
        Console} and select the \inlinecode{sa-brazil-1} region in the top-right
        corner.

  \item Open the service catalog and select \textbf{Flexus X}
        (under \textbf{Compute}).

  \item Click \textbf{Create Flexus X} to open the creation wizard.

  \item Under \textbf{Billing Mode}, select \textbf{Pay-per-use}.

  \begin{tip}
  Pay-per-use bills by the hour with no long-term commitment. You only pay
  for the time the instance is running. This is ideal for development and
  testing --- stop the instance when not in use to reduce costs.
  \end{tip}

  \item Set the compute resources using the slider:
        drag to \textbf{4 vCPU} and \textbf{16 GB RAM}.

  \begin{infobox}
  Flexus X uses a slider for vCPU and memory instead of fixed flavor names.
  The resulting flavor ID is \inlinecode{x1.4u.16g}, but you select it by
  dragging the vCPU and RAM controls --- there is no flavor dropdown.
  \end{infobox}

  \item Fill in the remaining basic configuration:

  \begin{code}
AZ:             Any (choose the availability zone you prefer)
Image:          Ubuntu 24.04 Server 64bit
Disk:           100 GB General Purpose SSD
VPC:            vpc-default
Subnet:         subnet-default
Security Group: sg-xgate (port 4444, Huawei internal IPs only)
  \end{code}

  \begin{warning}
  The EIP is auto-assigned and configured to be released when the instance is
  deleted. If you need a fixed IP that persists across stop/start cycles,
  allocate a dedicated EIP and bind it manually after creation.
  \end{warning}

  \item Under \textbf{EIP} settings, configure the public IP:

  \begin{code}
EIP:            Auto-assign
Billing:        Pay by traffic
Bandwidth:      1000 Mbit/s
Release:        Yes — release EIP when instance is deleted
  \end{code}

  \item Under \textbf{Cloud Eye}, enable monitoring (recommended).

  \begin{infobox}
  Cloud Eye is a monitoring service that provides functions such as real-time
  monitoring, alarm reporting, resource grouping, and website monitoring.
  Enabling it gives you visibility into instance health, CPU, memory, and
  network metrics out of the box.
  \end{infobox}

  \item Set the instance name to \inlinecode{flexusx-dev}
        (or any name you prefer).

  \item Under \textbf{Login Configuration}, select \textbf{Key pair} as the
        authentication method. If you do not yet have a key pair, click
        \textbf{Create Key Pair} to generate one inside Huawei Cloud and
        download the \inlinecode{.pem} file.

  \begin{infobox}
  The key pair is the primary access method. Later, you can set a root
  password on the instance (via the Console or after first login with the key)
  to have a second way to access the instance.
  \end{infobox}

  \item Click \textbf{Create Now} and wait for the status to change to
        \textbf{Running}.
\end{enumerate}

\subsection{Verify the instance}

On the instance details page, copy the \textbf{EIP} field. This is the address you
will use to connect in the next chapters.

\begin{table}[H]
  \begin{hutable}{|l|l|l|}
    \rowcolor{huaweired} \thd{Step} & \thd{What to check} & \thd{Expected result} \\
    \tbody
    Security group & \inlinecode{sg-xgate} with TCP 4444 rules   & 3 inbound rules \\
    Billing mode   & Pay-per-use selected                        & Hourly billing \\
    Compute        & 4 vCPU, 16 GB RAM via slider                & Flavor \inlinecode{x1.4u.16g} \\
    Image          & Ubuntu 24.04 Server 64bit                   & Selected \\
    EIP            & Auto-assigned                                & Public IP available \\
    Cloud Eye      & Monitoring enabled                           & Metrics visible \\
    Instance name  & \inlinecode{flexusx-dev}                     & Set \\
    Key pair       & Created and downloaded                       & \inlinecode{.pem} file saved \\
    Instance       & Status is \textbf{Running}                   & EIP copied \\
  \end{hutable}
  \caption{Chapter 1 verification checklist.}
\end{table}

% =====================================================================
%  Chapter 2 — Configure Flexus X for Remote Access
% =====================================================================
\section{Configure Flexus X for Remote Access}

\begin{objectives}
  \generalobjective{Prepare the Flexus X instance for remote access by setting a root
  password via the Console and changing the SSH port from 22 to 4444. This
  is required because the Huawei corporate proxy blocks outbound connections
  on port 22.}

  \prerequisites
  \begin{itemize}
    \item Chapter 1 completed — Flexus X instance running with a public IP.
    \item Access to the Huawei Cloud Console.
  \end{itemize}
\end{objectives}

\subsection{Set the root password}

\objective{Use the Console to reset the root password of the Flexus X instance.}

\stepbystep
\begin{enumerate}
  \item In the Console, navigate to \textbf{Flexus X} and click
        on the instance \inlinecode{flexusx-dev}.
  \item In the instance details page, click \menu{More, Reset Password}.

  \imagecap[width=0.75\linewidth]{assets/Flexus-Reset-Password.PNG}{Reset Password option in the
    Flexus X details page.}

  \item Enter a strong root password and confirm. Click \textbf{OK}.
  \item The instance will restart automatically. Wait for the status to return
        to \textbf{Running}.
\end{enumerate}

\begin{warning}
Choose a strong password (12+ characters, mixed case, digits, symbols).
This password enables remote login directly from the Console and serves as
a second access method alongside the SSH key pair.
\end{warning}

\subsection{Change the SSH port to 4444}

\objective{Log in via the Console remote login, edit \inlinecode{sshd\_config},
and change the SSH port from 22 to 4444.}

\stepbystep
\begin{enumerate}
  \item In the instance details page, click \textbf{Remote Login}.
        A browser-based terminal opens.

  \imagecap[width=0.75\linewidth]{assets/Flexus-Remote-Login.PNG}{Remote Login button in the Flexus X
    details page.}
  \item Log in as \inlinecode{root} using the password set in the previous
        subsection.
  \item Open the SSH daemon configuration file:

  \begin{code}[bash]
sudo vim /etc/ssh/sshd_config
  \end{code}

  \item Find the line \inlinecode{Port 22} (or \inlinecode{\#Port 22}) and change it
        to:

  \begin{code}
Port 4444
  \end{code}

  \item Save and exit vim (\inlinecode{:wq}).
  \item Reboot the instance from the Console to apply the change.
        The SSH daemon will start automatically on port 4444 after reboot.

  \imagecap[width=0.75\linewidth]{assets/Flexus-Restart.PNG}{Restart the Flexus X
    instance from the Console.}

  \begin{warning}
  The reboot is \textbf{mandatory} --- not optional. Until the instance is
  rebooted, SSH continues listening on port 22 and the proxy
  \inlinecode{CONNECT} to port 4444 will fail with
  \inlinecode{Ncat: Error reading proxy response Status-Line}. If you skip
  this step, all subsequent SSH connections through the proxy will fail.
  \end{warning}
\end{enumerate}

\begin{infobox}
After this change, SSH connections must use \inlinecode{Port 4444} explicitly.
The security group \inlinecode{sg-xgate} (created in Chapter 1) already allows
inbound TCP 4444 from the authorized Huawei internal IPs.
\end{infobox}

\subsection{Verify the configuration}

Use the table below to confirm all steps were completed successfully before
proceeding to the next chapter:

\begin{table}[H]
  \begin{hutable}{|l|l|l|}
    \rowcolor{huaweired} \thd{Step} & \thd{What to check} & \thd{Expected result} \\
    \tbody
    Security group & \inlinecode{sg-xgate} exists with TCP 4444 rules & 3 inbound rules \\
    Flexus X running & Instance \inlinecode{flexusx-dev} status        & \textbf{Running} \\
    Root password  & Remote Login works with password            & Shell prompt \\
    SSH port       & \inlinecode{sshd\_config} has \inlinecode{Port 4444}    & Config saved \\
    SSH restart    & Instance rebooted from Console              & Service active \\
  \end{hutable}
  \caption{Chapter 2 verification checklist.}
\end{table}

% =====================================================================
%  Chapter 3 — Install ncat and SSH Connection Topology
% =====================================================================
\section{Install ncat and SSH Connection Topology}

\begin{objectives}
  \generalobjective{Install \inlinecode{ncat} (from Nmap) on your local Windows
  machine and understand the SSH connection topology through the Huawei corporate
  proxy (xgate). ncat serves as the SSH \inlinecode{ProxyCommand}, tunneling SSH
  through the HTTP proxy via the \inlinecode{CONNECT} method.}

  \prerequisites
  \begin{itemize}
    \item Chapter 2 completed — Flexus X configured with root password and SSH
          on port 4444.
    \item The private key file (\inlinecode{.pem}) downloaded from Huawei Cloud.
  \end{itemize}
\end{objectives}

\begin{infobox}
\textbf{Premises:}
\begin{itemize}
  \item The proxy is \textbf{xgate} (Huawei corporate proxy) at
        \inlinecode{proxy.huawei.com:8080}.
  \item SSH runs on port \inlinecode{4444} (configured in Chapter 2).
  \item Login as \inlinecode{root} using the key pair from Chapter 1.
  \item \inlinecode{ncat} (from Nmap) is used as the \inlinecode{ProxyCommand}
        to tunnel SSH through the HTTP proxy via the \inlinecode{CONNECT} method.
  \item Proxy authentication is transparent --- it is handled by the corporate
        VPN, not by credentials in the SSH config.
\end{itemize}
\end{infobox}

\subsection{How the connection works}

When you specify a \inlinecode{ProxyCommand} in the SSH config, SSH launches
\inlinecode{ncat} as a child process. ncat performs the HTTP \inlinecode{CONNECT}
handshake with the proxy:

\begin{code}
ncat -> proxy:   CONNECT <EIP>:4444 HTTP/1.1
                 Host: <EIP>:4444

proxy -> ncat:   HTTP/1.1 200 Connection Established

[Both sides now send raw TCP bytes -- SSH negotiation begins over the tunnel]
\end{code}

After the \inlinecode{200 Connection Established} response, the proxy becomes a
blind TCP relay --- it neither inspects nor modifies the SSH traffic. SSH key
exchange, encryption, and authentication all happen inside this tunnel as if it
were a direct TCP connection.

\begin{infobox}
The corporate proxy blocks \inlinecode{CONNECT} to port 22 (standard SSH) but
allows it to port 4444. Enterprise HTTP proxies implement a \inlinecode{CONNECT}
port allowlist: port 443 (HTTPS) is always allowed, port 22 is usually blocked,
and port 4444 is often allowed because it is commonly used for HTTPS-alt. This
is why we configured SSH on port 4444 in Chapter 2.
\end{infobox}

\subsection{Why ncat}

\inlinecode{ncat} is the chosen proxy tunnel tool for several reasons:

\begin{itemize}
  \item \textbf{Full HTTP CONNECT support} --- ncat implements the
        \inlinecode{CONNECT} method natively via \inlinecode{--proxy-type http}.
  \item \textbf{Actively maintained} --- part of the Nmap project, with regular
        releases and bug fixes.
  \item \textbf{IPv4 forcing} --- the \inlinecode{-4} flag prevents failed IPv6
        attempts through the proxy.
  \item \textbf{Connect timeout} --- the \inlinecode{-w} flag prevents indefinite
        hangs if the proxy is unreachable.
  \item \textbf{Proxy auth support} --- \inlinecode{--proxy-auth} and the
        \inlinecode{NCAT\_PROXY\_AUTH} environment variable, if explicit
        credentials are ever needed.
\end{itemize}

\begin{infobox}
\inlinecode{connect.exe} (from Git for Windows) is a lighter alternative that
ships with Git, but it lacks the \inlinecode{-4} (force IPv4) and \inlinecode{-w}
(connect timeout) flags, which are important for reliable proxied connections.
\inlinecode{ncat} provides better control over the tunnel behavior.
\end{infobox}

\subsection{Install ncat on Windows}

\objective{Install Nmap, which provides \inlinecode{ncat.exe} --- the proxy
tunnel tool used by SSH \inlinecode{ProxyCommand}.}

\stepbystep
\begin{enumerate}
  \item Download the Nmap installer from
        \weblink{https://nmap.org/download.html\#windows}{nmap.org/download.html}.
        Use the \textbf{Latest stable release self-installer}
        (\inlinecode{nmap-7.991-setup.exe}, tested and verified).

  \item Run the installer (right-click \textbf{Run as Administrator}).

  \item During installation, uncheck \textbf{Npcap} --- it is a kernel driver
        for packet capture and is not needed for proxy tunneling. Also uncheck
        \textbf{Zenmap}, \textbf{Nping}, and \textbf{Ndiff} for a minimal install.

  \imagecap[width=0.75\linewidth]{assets/Ncat-Installation.PNG}{Nmap installer
    component selection --- uncheck Npcap, Zenmap, Nping, and Ndiff.}

  \begin{tip}
  Skipping Npcap avoids a kernel driver install and potential reboot.
  \inlinecode{ncat} in proxy mode uses standard Winsock2 TCP sockets and does
  not need Npcap at all.
  \end{tip}

  \item Choose the default install path:
        \inlinecode{C:\textbackslash Program Files (x86)\textbackslash Nmap}.

  \item After installation, verify that \inlinecode{ncat.exe} is available:

  \begin{code}[bash]
"C:\Program Files (x86)\Nmap\ncat.exe" --version
  \end{code}

  \item If the command shows the version string, the installation is correct.
\end{enumerate}

\begin{warning}
Nmap is commonly flagged by antivirus software due to port-scanning heuristics.
Windows Defender may quarantine \inlinecode{ncat.exe} on download or block
execution. Add an exclusion for the Nmap folder and the \inlinecode{ncat.exe}
process in \menu{Windows Security, Virus \& threat protection, Exclusions}.
If \inlinecode{ncat.exe} was already quarantined,
restore it from Protection history before adding the exclusion.
\end{warning}

\begin{table}[H]
  \begin{hutable}{|l|l|l|}
    \rowcolor{huaweired} \thd{Step} & \thd{What to check} & \thd{Expected result} \\
    \tbody
    Nmap download   & \inlinecode{nmap-7.991-setup.exe} saved      & Installer ready \\
    Components      & Npcap, Zenmap, Nping, Ndiff unchecked        & Minimal install \\
    Install path    & \inlinecode{C:\textbackslash Program Files (x86)\textbackslash Nmap} & Default path \\
    ncat.exe        & \inlinecode{ncat --version} shows version    & Installed correctly \\
    Antivirus       & Exclusion added for Nmap folder              & \inlinecode{ncat.exe} not blocked \\
  \end{hutable}
  \caption{Chapter 3 verification checklist.}
\end{table}

% =====================================================================
%  Chapter 4 — Install VS Code and Connect via Remote-SSH
% =====================================================================
\section{Install VS Code and Connect via Remote-SSH}

\begin{objectives}
  \generalobjective{Install VS Code from scratch on your local Windows machine,
  install the Remote-SSH extension, configure the SSH config file using VS Code
  as the editor, and connect to the Flexus X instance.}

  \prerequisites
  \begin{itemize}
    \item Chapter 3 completed --- ncat installed and proxy tunnel understood.
    \item The private key file (\inlinecode{.pem}) downloaded from Huawei Cloud.
  \end{itemize}
\end{objectives}

\subsection{Install VS Code}

\stepbystep
\begin{enumerate}
  \item Download VS Code from
        \weblink{https://code.visualstudio.com/download}{code.visualstudio.com/download}.
        Choose the \textbf{Windows} installer (User Installer, 64-bit).

  \item Run the installer with the default settings.

  \item After installation, verify VS Code is available:

  \begin{code}[bash]
code --version
  \end{code}
\end{enumerate}

\subsection{Install the Remote-SSH extension}

\stepbystep
\begin{enumerate}
  \item Open VS Code on your local Windows machine.
  \item Install the \textbf{Remote - SSH} extension (publisher: Microsoft):

  \begin{code}[bash]
code --install-extension ms-vscode-remote.remote-ssh
  \end{code}
\end{enumerate}

\subsection{Configure the SSH config file}

\objective{Create an SSH config entry that tunnels through the Huawei proxy
using \inlinecode{ncat} from Nmap. Use VS Code as the editor.}

\stepbystep
\begin{enumerate}
  \item Open or create the SSH config file using VS Code:

  \begin{code}[bash]
code %USERPROFILE%\.ssh\config
  \end{code}

  \item Add the following configuration. Replace these placeholders with your
        values:
        \begin{itemize}
          \item \inlinecode{<host-name>} --- a name of your choice (e.g., \inlinecode{flexusx-dev})
          \item \inlinecode{<EIP>} --- the Flexus X public IP copied from Chapter 1
          \item \inlinecode{<user>} --- your Windows username (e.g., \inlinecode{w00946241})
          \item \inlinecode{<key-file>} --- the name of your \inlinecode{.pem} key file
                downloaded from Huawei Cloud (e.g., \inlinecode{KeyPair-opencode})
        \end{itemize}

  \begin{code}
Host <host-name>
    HostName <EIP>
    User root
    Port 4444
    IdentityFile "C:\Users\<user>\.ssh\<key-file>.pem"
    ProxyCommand "C:\Program Files (x86)\Nmap\ncat.exe" -4 -w 10 --proxy proxy.huawei.com:8080 --proxy-type http %h %p
    ServerAliveInterval 15
    ServerAliveCountMax 4
    ConnectTimeout 10
    TCPKeepAlive yes
    StrictHostKeyChecking accept-new
  \end{code}

  \begin{infobox}
  The \inlinecode{IdentityFile} path has two placeholders:
  \inlinecode{<user>} is your Windows username (the folder under
  \inlinecode{C:\textbackslash Users\textbackslash}), and \inlinecode{<key-file>}
  is the name of the \inlinecode{.pem} file downloaded from Huawei Cloud
  in Chapter 1. For example:
  \inlinecode{C:\textbackslash Users\textbackslash w00946241\textbackslash .ssh\textbackslash KeyPair-opencode.pem}.
  \end{infobox}
\end{enumerate}

\subsection{What each ncat flag does}

\begin{table}[H]
  \begin{hutable}{|l|l|}
    \rowcolor{huaweired} \thd{Flag} & \thd{Purpose} \\
    \tbody
    \inlinecode{-4} & Force IPv4 (avoids failed IPv6 attempts through the proxy) \\
    \inlinecode{-w 10} & 10-second connect timeout (prevents indefinite hang) \\
    \inlinecode{--proxy} & Proxy address and port (\inlinecode{proxy.huawei.com:8080}) \\
    \inlinecode{--proxy-type http} & Use HTTP CONNECT method \\
    \inlinecode{\%h} & Substituted by SSH with the HostName (the EIP) \\
    \inlinecode{\%p} & Substituted by SSH with the Port (4444) \\
  \end{hutable}
  \caption{ncat ProxyCommand flags.}
\end{table}

\begin{tip}
The \inlinecode{ServerAliveInterval 15} and \inlinecode{ServerAliveCountMax 4}
settings send a keepalive every 15 seconds and disconnect after 4 missed
keepalives (60 seconds dead detection). This is critical for proxied
connections, which may drop idle sessions after 5--30 minutes.
\end{tip}

\subsection{Connect to the Flexus X via VS Code}

\stepbystep
\begin{enumerate}
  \item Press \inlinecode{F1} to open the Command Palette, type
        \textbf{Remote-SSH: Connect to Host...}, and press \inlinecode{Enter}.

  \imagecap[width=0.75\linewidth]{assets/SSH-Connection-1.PNG}{VS Code Command
    Palette --- select Remote-SSH: Connect to Host...}

  \item Select your host name (e.g., \inlinecode{flexusx-dev}) from the list
        of configured hosts.

  \item If prompted to select the platform type, choose \textbf{Linux}.
        This tells VS Code which server binary to download and avoids
        auto-detection delays.

  \begin{infobox}
  VS Code only asks for the platform type the first time you connect to a
  new host. The selection is saved in your local settings
  (\inlinecode{remote.SSH.remotePlatform}) and reused on subsequent connections.
  \end{infobox}

  \item VS Code opens a new window and begins connecting. On the first
        connection, it downloads and installs the VS Code Server on the
        remote instance --- this takes about 1--2 minutes. Subsequent
        connections are faster.

  \item Once connected, click \textbf{Open Folder} and select
        \inlinecode{/home} (or your working directory).

  \item Verify the connection by opening a terminal in VS Code
        (\inlinecode{Ctrl+`}) and running:

  \begin{code}[bash]
uname -a
df -h
  \end{code}
\end{enumerate}

\begin{infobox}
The first connection installs the VS Code Server on the remote instance, which
requires outbound internet access on port 443 (HTTPS). The Flexus X instance
has an EIP, so it typically has outbound internet access. If the connection
hangs during server installation, verify with
\inlinecode{curl -sI https://update.code.visualstudio.com} from the instance.
\end{infobox}

\subsection{Troubleshooting}

If the connection fails, check these common issues:

\begin{itemize}
  \item \textbf{Ncat: Error reading proxy response} --- SSH was not restarted
        after the port change (Chapter 2). \textbf{Fix:} reboot the instance
        from the Console.
  \item \textbf{Ncat: Error reading proxy response} --- proxy
        \inlinecode{CONNECT} permission missing for this IP.
        \textbf{Fix:} request IT permission for the instance EIP on port 4444.
  \item \textbf{Connection timed out} --- proxy unreachable or
        \inlinecode{CONNECT} blocked. \textbf{Fix:} verify proxy address; test
        ncat with the \inlinecode{-v} flag.
  \item \textbf{UNPROTECTED PRIVATE KEY FILE} --- key permissions too open.
        \textbf{Fix:} run \inlinecode{icacls} to restrict the key file to your
        user only.
  \item \textbf{VS Code hangs on first connect} --- instance cannot reach
        internet on port 443. \textbf{Fix:} verify
        \inlinecode{curl -sI https://update.code.visualstudio.com} from the
        instance.
  \item \textbf{\$PLATFORM is undefined} --- SSH connection failed before
        server install. \textbf{Fix:} check the Remote-SSH output log; fix the
        underlying SSH error first.
\end{itemize}

\begin{infobox}
To debug the proxy connection, run ncat directly with verbose output:
\begin{code}
"C:\Program Files (x86)\Nmap\ncat.exe" -v --proxy proxy.huawei.com:8080 --proxy-type http <EIP> 4444
\end{code}
This shows the \inlinecode{CONNECT} request and the proxy's raw response (or
connection close). For SSH-level debugging, run
\inlinecode{ssh -vvv <host-name> "echo ok"}.
\end{infobox}

\begin{table}[H]
  \begin{hutable}{|l|l|l|}
    \rowcolor{huaweired} \thd{Step} & \thd{What to check} & \thd{Expected result} \\
    \tbody
    VS Code         & \inlinecode{code --version} works            & Installed \\
    Remote-SSH      & Extension installed                          & Active \\
    SSH config      & Host entry with ncat ProxyCommand            & Config saved \\
    Connection      & Remote-SSH connects to host                  & Flexus X accessible \\
    Terminal        & \inlinecode{uname -a} and \inlinecode{df -h} & Output verified \\
  \end{hutable}
  \caption{Chapter 4 verification checklist.}
\end{table}

% =====================================================================
%  Chapter 5 — Install oh-my-coding-maas-gateway
% =====================================================================
\section{Install the MaaS Coding Gateway}

\begin{objectives}
  \generalobjective{Install the oh-my-coding-maas-gateway on the Flexus X
  instance. This is a self-hosted LiteLLM proxy that routes Huawei ModelArts MaaS
  models to coding tools with load balancing, virtual keys, dual-format API
  endpoints, and Prometheus + Grafana observability.}

  \prerequisites
  \begin{itemize}
    \item Chapter 4 completed — VS Code connected to the Flexus X via Remote-SSH.
    \item Docker and Docker Compose available (the installer handles this).
    \item A Huawei ModelArts MaaS API key.
  \end{itemize}
\end{objectives}

\subsection{Overview}

The oh-my-coding-maas-gateway is a self-hosted proxy gateway that sits between
your coding tools and Huawei ModelArts MaaS (Model-as-a-Service) models. It
provides:

\begin{itemize}
  \item \textbf{Load balancing} --- distributes requests across multiple MaaS
        API keys (N keys $\to$ N deployments per model per format).
  \item \textbf{Virtual key management} --- each coding tool gets its own
        virtual key with unlimited budget, isolated from the real MaaS keys.
  \item \textbf{Dual-format API endpoints} --- OpenAI Chat Completions
        (\inlinecode{/v1/chat/completions}) and Anthropic Messages
        (\inlinecode{/v1/messages}), so tools using different API formats can
        all reach the same MaaS models.
  \item \textbf{Budget tracking} --- per-key spend tracking via LiteLLM +
        PostgreSQL, with cost-per-token configured for all models.
  \item \textbf{Full observability} --- Prometheus metrics + Grafana 39-panel
        dashboard (latency, errors, throughput, tokens, cache, cost).
\end{itemize}

\begin{infobox}
The gateway repository is at
\weblink{https://github.com/wallacelw/oh-my-coding-maas-gateway}{wallacelw/oh-my-coding-maas-gateway}.
It supports 4 models: \inlinecode{glm-5.2} (1M context),
\inlinecode{glm-5.1}, \inlinecode{deepseek-v4-pro}, and
\inlinecode{deepseek-v4-flash} (384K max output). This guide was tested with
v1.10.9; the installation process is the same for newer versions.
\end{infobox}

\subsection{Architecture}

The gateway deploys 4 Docker containers via Docker Compose:

\begin{code}
  Tools               LiteLLM (:4000)                  Huawei MaaS
  ─────               ───────────────                  ────────────

  opencode    --> /v1/chat/completions --> openai/    --> MaaS OpenAI endpoint
  Codex CLI   --> /v1/responses        --> openai/    --> MaaS OpenAI endpoint
  Claude Code --> /v1/messages         --> anthropic/ --> MaaS Anthropic endpoint
  Pi agent    --> /v1/chat/completions --> openai/    --> MaaS OpenAI endpoint

  LiteLLM: load-balances across N MaaS keys . PostgreSQL (:5432)
  Observability: LiteLLM --/metrics--> Prometheus (:9090) --> Grafana (:3000)
\end{code}

\begin{table}[H]
  \begin{hutable}{|l|l|l|}
    \rowcolor{huaweired} \thd{Service} & \thd{Port} & \thd{Purpose} \\
    \tbody
    LiteLLM Proxy  & 4000 & API gateway (OpenAI + Anthropic format) \\
    PostgreSQL     & 5432 & LiteLLM database (keys, spend, config) \\
    Prometheus     & 9090 & Metrics storage and querying \\
    Grafana        & 3000 & Dashboard visualization (39 panels) \\
  \end{hutable}
  \caption{Docker services and ports.}
\end{table}

\begin{infobox}
All ports bind to \inlinecode{127.0.0.1} (localhost only) by default. To access
Grafana and the LiteLLM Admin UI from your local machine, use SSH port
forwarding in VS Code or add port forwards in the Remote-SSH connection.
\end{infobox}

Each model has two deployments --- one OpenAI, one Anthropic --- so all tools
can use the same MaaS models regardless of their API format. The
\inlinecode{claude-} prefix (e.g., \inlinecode{claude-glm-5.2}) prevents routing
conflicts in LiteLLM.

\subsection{Run the bootstrap installer}

\objective{Install the gateway using the one-command bootstrap script.}

\stepbystep
\begin{enumerate}
  \item In the VS Code remote terminal, run the bootstrap script:

  \begin{code}[bash]
curl -fsSL https://raw.githubusercontent.com/wallacelw/oh-my-coding-maas-gateway/main/scripts/bootstrap.sh | bash
  \end{code}

  \item The installer will:
    \begin{itemize}
      \item Prompt for your Huawei MaaS API key.
      \item Show an interactive menu to select which coding tools to install
            (opencode, Codex CLI, Claude Code, Pi agent).
      \item Auto-install prerequisites (Docker, Node.js, etc.).
      \item Deploy the proxy + observability stack (LiteLLM, PostgreSQL,
            Prometheus, Grafana).
      \item Configure each tool with its own virtual key.
      \item Run 120 validation checks and offer to install the companion skill.
    \end{itemize}

  \begin{warning}
  Estimated time: \textasciitilde5 min for a fresh install, \textasciitilde2 min
  for an upgrade. If you are behind a proxy, ensure \inlinecode{http\_proxy} and
  \inlinecode{https\_proxy} environment variables are set before running the script.
  \end{warning}
\end{enumerate}

\begin{tip}
For non-interactive installs (CI/automation), pass the API key directly. The
\inlinecode{-y} flag suppresses all prompts (prerequisites, tool selection,
skill install):
\begin{code}
curl -fsSL https://raw.githubusercontent.com/wallacelw/oh-my-coding-maas-gateway/main/scripts/bootstrap.sh | bash -s -- -y --api-key=sk-xxxx
\end{code}
\end{tip}

\subsection{Verify the installation}

\stepbystep
\begin{enumerate}
  \item Check that the proxy is running:

  \begin{code}[bash]
curl http://localhost:4000/health
  \end{code}

  \item Open the Grafana dashboard at \inlinecode{http://localhost:3000} to
        monitor usage and metrics (39 panels across 7 sections).

  \item Open the LiteLLM Admin UI at \inlinecode{http://localhost:4000/ui} to
        manage virtual keys and model routing.

  \item Launch your coding tool:

  \begin{code}[bash]
opencode          # or:  codex  or:  claude --bare  or:  pi
  \end{code}
\end{enumerate}

\begin{tip}
To upgrade the gateway later, re-run the same bootstrap command. It detects
existing installs, compares versions, and pulls updates --- all secrets and
data are preserved.
\end{tip}

\begin{table}[H]
  \begin{hutable}{|l|l|l|}
    \rowcolor{huaweired} \thd{Step} & \thd{What to check} & \thd{Expected result} \\
    \tbody
    Bootstrap       & Installer completed                          & All services deployed \\
    Proxy health    & \inlinecode{curl localhost:4000/health}     & Healthy \\
    Grafana         & \inlinecode{localhost:3000}                  & Dashboard accessible \\
    LiteLLM UI      & \inlinecode{localhost:4000/ui}               & Admin UI accessible \\
    Coding tool     & Launched and connected                       & Gateway routing \\
  \end{hutable}
  \caption{Chapter 5 verification checklist.}
\end{table}

% =====================================================================
%  Chapter 6 — Install the Huawei Project
% =====================================================================
\section{Install the Huawei Document Templates Project}

\begin{objectives}
  \generalobjective{Clone the Huawei document templates repository on the Flexus X
  instance, run the setup script, and verify that the LaTeX toolchain and
  the coding agent skill are ready for use. The script assumes a fresh
  Ubuntu install --- it installs all dependencies from scratch and
  configures VS Code at the user level.}

  \prerequisites
  \begin{itemize}
    \item Chapter 5 completed — MaaS gateway installed and running.
    \item Git installed (\inlinecode{sudo apt install git}).
    \item The repository URL for the Huawei document templates project.
    \item A fresh Ubuntu 24.04 instance (no prior LaTeX installation
          required).
  \end{itemize}
\end{objectives}

\begin{infobox}
The \inlinecode{install.sh} script is idempotent --- it is safe to run multiple
times. It merges VS Code settings rather than overwriting them, and
\inlinecode{apt-get install} skips packages that are already present.
\end{infobox}

\subsection{Clone and set up the project}

\stepbystep
\begin{enumerate}
  \item Clone the repository:

  \begin{code}[bash]
cd /home
git clone https://github.com/wallacelw/huawei-doc-templates.git
cd huawei-doc-templates
  \end{code}

  \item Run the setup script (installs XeLaTeX, latexmk, fonts, the coding
        agent skill, and configures VS Code LaTeX Workshop at the user
        level):

  \begin{code}[bash]
./install.sh
  \end{code}

  \begin{warning}
  The script runs \inlinecode{sudo apt-get install} for TeX Live packages. If you
  do not have sudo access, ask your administrator to install the packages
  listed in \inlinecode{install.sh}.
  \end{warning}

  \item Verify the LaTeX toolchain:

  \begin{code}[bash]
xelatex --version
latexmk --version
  \end{code}
\end{enumerate}

\subsection{Compile the sample guides}

The project includes a self-documenting Makefile. To compile all samples
(guide + technical, Portuguese + English):

\stepbystep
\begin{enumerate}
  \item From the repository root, run:

  \begin{code}[bash]
make samples
  \end{code}

  \item Verify the PDFs were generated:

  \begin{code}[bash]
ls -l examples/guide/pt/main.pdf
ls -l examples/guide/en/main.pdf
ls -l examples/technical/pt/main.pdf
ls -l examples/technical/en/main.pdf
  \end{code}
\end{enumerate}

\begin{tip}
To compile a single project, use \inlinecode{make project DIR=examples/guide/pt}
or run \inlinecode{latexmk main.tex} from the project's \inlinecode{src/}
directory. Run \inlinecode{make} with no arguments to see all available targets.
\end{tip}

\subsection{Create your first guide with the skill}

\stepbystep
\begin{enumerate}
  \item In the VS Code remote terminal, launch your coding agent:

  \begin{code}[bash]
cd /home/huawei-doc-templates
opencode          # or: codex, claude --bare, pi
  \end{code}

  \item Run the guide skill:

  \begin{code}
/skill huawei-template-guide
  \end{code}

  \item The skill will ask for:
    \begin{itemize}
      \item \textbf{Title} — e.g. ``Provisioning an OBS Bucket''
      \item \textbf{Language} — Portuguese or English
      \item \textbf{Project name} — used as the folder name
      \item \textbf{Location} — where to create the project folder
    \end{itemize}

  \item The skill generates the \inlinecode{.tex} file, creates a
        \inlinecode{.latexmkrc} with the correct \inlinecode{TEXINPUTS}, compiles the
        document, and reports the page count.

  \badge{Ready}
\end{enumerate}

\begin{infobox}
The skill is auto-discovered via \inlinecode{opencode.json} which registers
\inlinecode{templates/} as a skill discovery path. No manual skill installation
is needed when working inside the repo.
\end{infobox}

\subsection{Write content using template commands}

The template provides commands for common guide elements:

Headings (automatic numbering, H1 starts on a new page):

\begin{code}
\section{Chapter Title}          % H1: giant number + title + rule
\subsection{Section Title}        % H2
\subsubsection{Subsection Title}  % H3
\end{code}

Objectives block (closes with a horizontal rule):

\begin{code}
\begin{objectives}
  \generalobjective{...}
  \prerequisites
  \begin{itemize}
    \item ...
  \end{itemize}
\end{objectives}
\end{code}

Step-by-step and objectives (used inside subsections):

\begin{code}
\objective{...}
\stepbystep
\begin{enumerate}
  \item ...
\end{enumerate}
\end{code}

Callout boxes (environment syntax):

\begin{itemize}
  \item \inlinecode{warning} --- amber warning box
  \item \inlinecode{tip} --- green tip box
  \item \inlinecode{infobox} --- blue info box
\end{itemize}

Code blocks and inline commands:

\begin{itemize}
  \item \inlinecode{\textbackslash begin\{code\}} --- code block (verbatim,
        no escaping needed inside)
  \item \inlinecode{\textbackslash begin\{code\}[bash]} --- same, with a
        language hint
  \item \inlinecode{\textbackslash inlinecode\{...\}} --- inline monospace code
\end{itemize}

Images, links, and other inline commands:

\begin{code}
\image{assets/screenshot.png}        % centered image
\imagecap{assets/screenshot.png}{caption}  % image with numbered caption
\note{Italic observation.}            % note paragraph
\weblink{https://...}{link text}     % blue clickable link
\menu{VPC, Security Groups}          % menu path: VPC -> Security Groups
\badge{New}                          % inline red label
\end{code}

\begin{infobox}
The project also includes a \inlinecode{technical} template for technical
reports. Use \inlinecode{/skill huawei-template-technical} to create technical
reports with a 5-section structure (problem, root cause analysis, root cause,
trigger condition, workaround). See \inlinecode{templates/technical/SKILL.md}
for the full command reference.
\end{infobox}

\begin{table}[H]
  \begin{hutable}{|l|l|l|}
    \rowcolor{huaweired} \thd{Step} & \thd{What to check} & \thd{Expected result} \\
    \tbody
    Repository      & Cloned to \inlinecode{/home/huawei-doc-templates} & Files present \\
    install.sh      & Completed without errors                     & Dependencies installed \\
    XeLaTeX         & \inlinecode{xelatex --version} works        & Toolchain ready \\
    latexmk         & \inlinecode{latexmk --version} works        & Build system ready \\
    Sample PDFs     & \inlinecode{make samples} compiles all       & Guide + technical, PT + EN \\
    Skill           & \inlinecode{/skill huawei-template-guide}    & Available \\
  \end{hutable}
  \caption{Chapter 6 verification checklist.}
\end{table}

% =====================================================================
%  Chapter 7 — Clean Up
% =====================================================================
\section{Clean Up}

\begin{objectives}
  \generalobjective{Remove the resources created in this guide to avoid ongoing
  charges on your Huawei Cloud account.}
\end{objectives}

\stepbystep
\begin{enumerate}
  \item In the Console, navigate to \textbf{Flexus X}.
  \item Select the instance \inlinecode{flexusx-dev}, click \textbf{More} and
        choose \textbf{Delete}.
  \item Confirm deletion. The EIP is released automatically (configured in
        Chapter 1 to release on instance deletion).
  \item Navigate to \menu{VPC, Security Groups},
        select \inlinecode{sg-xgate}, and delete it.
  \item Remove the SSH config entry for \inlinecode{flexusx-dev} from your
        Windows SSH config file
        (\inlinecode{C:\textbackslash Users\textbackslash <user>\textbackslash .ssh\textbackslash config}).
\end{enumerate}

\begin{warning}
If you used Terraform to provision the infrastructure, run
\inlinecode{terraform destroy} instead to remove all resources declared in the
configuration.
\end{warning}

\begin{table}[H]
  \begin{hutable}{|l|l|l|}
    \rowcolor{huaweired} \thd{Step} & \thd{What to check} & \thd{Expected result} \\
    \tbody
    Flexus X instance & \inlinecode{flexusx-dev} deleted            & No running instance \\
    EIP               & Released                                    & No public IP charges \\
    Security group    & \inlinecode{sg-xgate} deleted               & No SG charges \\
    SSH config        & Host entry removed                          & Clean config \\
  \end{hutable}
  \caption{Chapter 7 verification checklist.}
\end{table}

\note{This guide was generated using the \inlinecode{guide.cls} template. For the
full command reference, see the template's README.md or run
\inlinecode{/skill huawei-template-guide} in your coding agent.}

% =====================================================================
%  Changelog
% =====================================================================
\begin{changelog}
  \changelogentry{3.1.1}{\today}{\item Updated MaaS Gateway chapter to v1.10.9: non-interactive mode (\inlinecode{-y}) now suppresses all prompts (prerequisites, tool selection, skill install). \item Changed project clone path from \inlinecode{\textasciitilde} (root home) to \inlinecode{/home} as default working directory. \item Added project standards to AGENTS.md: workflow, end-to-end validation, code style, git conventions, when unsure.}
  \changelogentry{3.1.0}{\today}{\item Switched from ECS to Flexus X instance (4 vCPU, 16 GB RAM via slider, flavor \inlinecode{x1.4u.16g}) --- more cost-efficient and robust. Updated creation steps to match Flexus X console order: billing mode first, vCPU/RAM slider (no flavor dropdown), Cloud Eye monitoring (recommended), then instance name and key pair login. Replaced screenshots with Flexus X console images. \item Rewrote Chapter 3: switched from \inlinecode{connect.exe} to \inlinecode{ncat} (from Nmap) as the SSH ProxyCommand. Added explanation of HTTP CONNECT method, why port 4444 works, why ncat is chosen, ncat installation with antivirus exclusions, flag reference table, and improved keepalive settings. \item Restructured chapters: moved SSH config and VS Code connection to new Chapter 4 (Install VS Code and Connect via Remote-SSH). Chapters 5--7 renumbered. \item Added verification checklist tables to all 7 chapters. \item Added mandatory reboot warning in Chapter 2 (SSH port change). Added troubleshooting subsection in Chapter 4 with common SSH connection issues and ncat debugging commands.}
  \changelogentry{3.0.1}{\today}{\item Minor fixes: test-filter.sh auto-discovery, build.sh help text, Makefile aggregate descriptions, 5-section comment correction, CHANGELOG version gap note, callout-in-code fix, CI font auto-discovery, dofile path comment, test-sync.sh header.}
  \changelogentry{3.0.0}{\today}{\item Modular pipeline refactoring: shared Lua filter factory, shared DOCX fix logic, build system auto-discovery, unified format pipeline, GitHub Actions CI.}
  \changelogentry{2.16.0}{\today}{\item Updated to match template v2.16.0 — Lua filter fixes, build system updates, and technical template parity improvements.}
  \changelogentry{2.13.0}{2026-08-24}{
    \item Foundation refactoring Phase 4: make DOCX \texttt{--fix} post-processing robust. Replaced all silent-skip guards in \texttt{create-reference-docx.py} with loud \texttt{RuntimeError} assertions for expected XML elements (Heading1-4, Title, VerbatimChar, Normal, docDefaults, numbering indentation). Added pandoc version pin (3.1.0--3.2.0) to catch output-structure regressions on upgrade. Extended \texttt{test-docx-fix.sh} with version check and 4 loud-failure tests (missing style $\to$ non-zero exit).
  }
  \changelogentry{2.12.0}{2026-08-24}{
    \item Foundation refactoring Phase 3: introduce \texttt{parse\_preamble()} in the Lua filter, extracting all class options (\texttt{portuguese}, \texttt{indentbody}, \texttt{notime}, \texttt{nochangelog}) and all \texttt{\textbackslash set*} command values into a single preamble state object. Fixes 6 L18 divergences at root cause: English TOC label (``Table of Contents'' $\to$ ``Contents''), \texttt{[notime]} now gates time in DOCX/HTML, \texttt{[nochangelog]} now suppresses changelog in non-PDF formats, \texttt{\textbackslash setcoverlogo} honored instead of hardcoded, \texttt{\textbackslash setheaderlogo} warns if set in non-PDF, \texttt{\textbackslash note} renders as italic in all formats (matching PDF). Any future preamble command or class option now flows through automatically.
  }
  \changelogentry{2.11.0}{2026-08-24}{
    \item Foundation refactoring Phase 2: extend test matrix to cover all 3 output formats $\times$ 3 documents (9 cells). \texttt{round-trip.sh} now generates DOCX, counts structures in \texttt{word/document.xml}, and checks cross-format consistency (headings, images, code blocks, tables, callouts, 0 raw LaTeX across MD+DOCX+HTML). Added \texttt{test-docx-fix.sh} smoke test asserting H1 red border, heading colors, list indentation, and footer PAGE field in patched \texttt{styles.xml}.
  }
  \changelogentry{2.10.0}{2026-08-24}{
    \item Foundation refactoring Phase 1: sync cls version to project tag, fix stale doc references (L16 output count, L-range, setup-guide feature claim), add test-sync.sh for version consistency checking.
  }
  \changelogentry{2.9.1}{2026-08-23}{
    \item HTML and Markdown outputs now embed images as base64 data URIs. HTML uses pandoc \texttt{--embed-resources}, MD uses post-processing with \texttt{embed-images.py}. Outputs are self-contained — no external image files needed.
  }
  \changelogentry{2.9.0}{2026-08-23}{
    \item DOCX: fix content width (8504→9638tw, margins are 2cm not 3cm). Add footer page numbers. Style TOC heading (22pt bold + rule). Add hutable table width + fixed layout.
  }
  \changelogentry{2.8.3}{2026-08-23}{
    \item DOCX fix: callout tables now have explicit width (8504tw) and fixed layout. Changelog section heading now uses custom H1 format (56pt number, red border, right tab) instead of pandoc's default.
  }
  \changelogentry{2.8.2}{2026-08-23}{
    \item DOCX H1 heading: reverted table approach to paragraph with tab stop + red bottom border. Table broke navigation pane and TOC field. Paragraph preserves Heading1 style for navigation while matching PDF layout (number left, title right, red rule below).
  }
  \changelogentry{2.8.1}{2026-08-23}{
    \item DOCX heading fix: text color reverted to black (was incorrectly changed to Huawei red in v2.8.0). H1 now uses a two-column table (number left, title right) with red bottom rule. H2-H4 use simple inline format (number + title, no right tab stop).
  }
  \changelogentry{2.8.0}{2026-08-23}{
    \item DOCX output: match PDF styling --- heading colors (Huawei red), H1 number size (56pt), callout borders (left-only) and padding, changelog formatting (rules, bold version, italic date), caption style (9pt bold), badge style (8pt bold white on red), list indentation.
  }
  \changelogentry{2.7.1}{2026-08-17}{
    \item Rename project repository to all lowercase (\texttt{huawei-doc-templates}).
  }
  \changelogentry{2.7.0}{2026-08-12}{
    \item The \texttt{changelog} environment now emits its own section heading; \texttt{[nochangelog]} suppresses the heading and entries in one switch.
  }
  \changelogentry{2.6.1}{2026-08-12}{
    \item Fix \texttt{hutable} centering and row-color fill (switched to \texttt{\textbackslash hline} for \texttt{colortbl} compatibility).
  }
  \changelogentry{2.6.0}{2026-08-12}{
    \item Switched verification table to the new \texttt{hutable} full-grid style.
    \item Floats now default to \texttt{[H]} placement (in-source order).
  }
  \changelogentry{2.5.1}{2026-08-10}{
    \item H1 section title size reduced from 27pt to 20pt.
  }
  \changelogentry{2.5.0}{2026-08-10}{
    \item Default image size increased: width 65\% $\to$ 90\% of text width,
          height 40\% $\to$ 50\% of text height.
  }
  \changelogentry{2.4.0}{2026-08-10}{
    \item Enlarged H1 section title (18pt $\to$ 27pt, 50\% bigger).
    \item \texttt{nochangelog} option now hides version, date, and time on the
          cover page.
  }
  \changelogentry{2.3.1}{2026-08-09}{
    \item Rolled back \texttt{/ActualText} markers — broke code block rendering.
  }
  \changelogentry{2.3.0}{2026-08-09}{
    \item Added key-value sizing options to \inlinecode{\textbackslash image} and
          \inlinecode{\textbackslash imagecap} — \inlinecode{width} and
          \inlinecode{height} can now be set independently.
  }
  \changelogentry{2.2.4}{2026-08-09}{
    \item Fixed PDF copy-paste of URLs in code blocks — disabled Cascadia Code
          programming ligatures (\texttt{calt}) that corrupted \texttt{//} into
          \texttt{))} on copy.
  }
  \changelogentry{2.2.3}{2026-08-09}{
    \item Fixed italic font rendering — HarmonyOS Sans and Cascadia Code have
          no italic variants; enabled synthetic slant (\texttt{AutoFakeSlant}).
  }
  \changelogentry{2.2.2}{2026-08-08}{
    \item Fixed chapter reference in prerequisite (Chapter 3 → Chapter 4).
  }
  \changelogentry{2.2.1}{2026-08-08}{
    \item Changed table and figure caption labels to black (was Huawei red).
  }
  \changelogentry{2.2.0}{2026-08-08}{
    \item Added Huawei-branded table styling: red header bar, alternating row
          colors, and bordered cells.
  }
  \changelogentry{2.1.0}{2026-08-06}{
    \item Updated ECS flavor to \inlinecode{ac8.xlarge.2} (4 vCPU, 8 GB RAM).
    \item Updated OS image to Ubuntu 24.04 Server 64bit.
    \item Increased disk size to 100 GB.
    \item Recommended pay-per-use billing mode.
    \item EIP auto-assigned: pay by traffic, 1000 Mbit/s, release with ECS.
    \item Key pair authentication (created in Huawei Cloud); root password
          can be set later as a second access method.
    \item Security group \inlinecode{sg-xgate} now created \textbf{before} the ECS
          instance.
    \item Replaced SSH port 22 with port 4444 (Huawei proxy blocks port 22).
    \item Restricted SSH access to three specific Huawei internal IPs
          (\inlinecode{119.8.89.193}, \inlinecode{119.8.89.3}, \inlinecode{119.8.193.3}).
    \item Added \inlinecode{Port 4444} to SSH config examples.
  }
  \changelogentry{2.0.0}{2026-08-05}{
    \item Replaced OpenCode with oh-my-coding-maas-gateway as the default
          coding environment.
    \item Added changelog and versioning support to the template
          (\inlinecode{\textbackslash setdocversion}, \inlinecode{\textbackslash setdocdate},
          \inlinecode{changelog} environment).
    \item Renamed \inlinecode{info} environment to \inlinecode{infobox} to avoid
          package name collisions.
    \item Fixed external file code block implementation (\inlinecode{\textbackslash codefile}).
  }
  \changelogentry{1.0.0}{2026-08-05}{
    \item Initial version.
    \item Added ECS provisioning on Huawei Cloud.
    \item Added SSH with proxy configuration for VS Code Remote-SSH.
    \item Replaced OpenCode with oh-my-coding-maas-gateway (LiteLLM proxy
          with load balancing, virtual keys, and Grafana observability).
    \item Added Huawei document templates project setup.
    \item Added changelog and versioning support to the template.
  }
\end{changelog}

\end{document}
